%-----------------------------------------------------------------------------------------
% Autor dieser Vorlage:
% Stefan Macke (http://fachinformatiker-anwendungsentwicklung.net)
% Permalink zur Vorlage: http://fiae.link/LaTeXVorlageFIAE
%
% Sämtliche verwendeten Abbildungen, Tabellen und Listings stammen von Dirk Grashorn.
%
% Lizenz: Creative Commons 4.0 Namensnennung - Weitergabe unter gleichen Bedingungen
% -----------------------------------------------------------------------------------------

\documentclass[
	ngerman,
	toc=listof, % Abbildungsverzeichnis sowie Tabellenverzeichnis in das Inhaltsverzeichnis aufnehmen
	toc=bibliography, % Literaturverzeichnis in das Inhaltsverzeichnis aufnehmen
	footnotes=multiple, % Trennen von direkt aufeinander folgenden Fußnoten
	parskip=half, % vertikalen Abstand zwischen Absätzen verwenden anstatt horizontale Einrückung von Folgeabsätzen
	numbers=noendperiod % Den letzten Punkt nach einer Nummerierung entfernen (nach DIN 5008)
]{scrartcl}
\pdfminorversion=5 % erlaubt das Einfügen von pdf-Dateien bis Version 1.7, ohne eine Fehlermeldung zu werfen (keine Garantie für fehlerfreies Einbetten!)
\usepackage[utf8]{inputenc} % muss als erstes eingebunden werden, da Meta/Packages ggfs. Sonderzeichen enthalten

% !TEX root = Projektdokumentation.tex

% Hinweis: der Titel muss zum Inhalt des Projekts passen und den zentralen Inhalt des Projekts deutlich herausstellen
\newcommand{\titel}{Nuxt3 Typo3 Skeleton}
\newcommand{\untertitel}{Webbasiertes Headless Content-Management-System}
\newcommand{\kompletterTitel}{\titel{} -- \untertitel}

\newcommand{\autorName}{Lukas Röding}
\newcommand{\autorAnschrift}{Hohe Straße 8}
\newcommand{\autorOrt}{44139 Dortmund}

\newcommand{\betriebLogo}{Firmenlogo.png}
\newcommand{\betriebName}{bits & likes GmbH}
\newcommand{\betriebAnschrift}{Rheinische Str. 171}
\newcommand{\betriebOrt}{44139 Dortmund}

\newcommand{\ausbildungsberuf}{Fachinformatiker für Anwendungsentwicklung}
\newcommand{\betreff}{Dokumentation zur betrieblichen Projektarbeit}
\newcommand{\pruefungstermin}{Winter 2022/23}
\newcommand{\abgabeOrt}{Dortmund}
\newcommand{\abgabeTermin}{23.04.2015}
 % Metadaten zu diesem Dokument (Autor usw.)
% !TEX root = ../Projektdokumentation.tex

% Anpassung an Landessprache ---------------------------------------------------
\usepackage{babel}

% Umlaute ----------------------------------------------------------------------
%   Umlaute/Sonderzeichen wie äüöß direkt im Quelltext verwenden (CodePage).
%   Erlaubt automatische Trennung von Worten mit Umlauten.
% ------------------------------------------------------------------------------
\usepackage[T1]{fontenc}
\usepackage{textcomp} % Euro-Zeichen etc.

% Schrift ----------------------------------------------------------------------
\usepackage{lmodern} % bessere Fonts
\usepackage{relsize} % Schriftgröße relativ festlegen

% Tabellen ---------------------------------------------------------------------
\PassOptionsToPackage{table}{xcolor}
\usepackage{tabularx}
% für lange Tabellen
\usepackage{longtable}
\usepackage{array}
\usepackage{ragged2e}
\usepackage{lscape}
\newcolumntype{w}[1]{>{\raggedleft\hspace{0pt}}p{#1}} % Spaltendefinition rechtsbündig mit definierter Breite
\usepackage{float}
% Grafiken ---------------------------------------------------------------------
\usepackage[dvips,final]{graphicx} % Einbinden von JPG-Grafiken ermöglichen
\usepackage{graphics} % keepaspectratio
\usepackage{floatflt} % zum Umfließen von Bildern
\graphicspath{{Bilder/}} % hier liegen die Bilder des Dokuments

% Sonstiges --------------------------------------------------------------------
\usepackage[titles]{tocloft} % Inhaltsverzeichnis DIN 5008 gerecht einrücken
\usepackage{amsmath,amsfonts} % Befehle aus AMSTeX für mathematische Symbole
\usepackage{enumitem} % anpassbare Enumerates/Itemizes
\usepackage{xspace} % sorgt dafür, dass Leerzeichen hinter parameterlosen Makros nicht als Makroendezeichen interpretiert werden

\usepackage{makeidx} % für Index-Ausgabe mit \printindex
\usepackage[printonlyused]{acronym} % es werden nur benutzte Definitionen aufgelistet

% Einfache Definition der Zeilenabstände und Seitenränder etc.
\usepackage{setspace}
\usepackage{geometry}

% Symbolverzeichnis
\usepackage[intoc]{nomencl}
\let\abbrev\nomenclature
\renewcommand{\nomname}{Abkürzungsverzeichnis}
\setlength{\nomlabelwidth}{.25\hsize}
\renewcommand{\nomlabel}[1]{#1 \dotfill}
\setlength{\nomitemsep}{-\parsep}

\usepackage{varioref} % Elegantere Verweise. „auf der nächsten Seite“
\usepackage{url} % URL verlinken, lange URLs umbrechen etc.


\usepackage{chngcntr} % fortlaufendes Durchnummerieren der Fußnoten
% \usepackage[perpage]{footmisc} % Alternative: Nummerierung der Fußnoten auf jeder Seite neu

\usepackage{ifthen} % bei der Definition eigener Befehle benötigt
\usepackage{todonotes} % definiert u.a. die Befehle \todo und \listoftodos
\usepackage[square]{natbib} % wichtig für korrekte Zitierweise

% PDF-Optionen -----------------------------------------------------------------
\usepackage{pdfpages}
\pdfminorversion=5 % erlaubt das Einfügen von pdf-Dateien bis Version 1.7, ohne eine Fehlermeldung zu werfen (keine Garantie für fehlerfreies Einbetten!)
\usepackage[
    bookmarks,
    bookmarksnumbered,
    bookmarksopen=true,
    bookmarksopenlevel=1,
    colorlinks=true,
% diese Farbdefinitionen zeichnen Links im PDF farblich aus
    anchorcolor=AOBlau,% Ankertext
    citecolor=AOBlau, % Verweise auf Literaturverzeichniseinträge im Text
    filecolor=AOBlau, % Verknüpfungen, die lokale Dateien öffnen
    menucolor=AOBlau, % Acrobat-Menüpunkte
    urlcolor=AOBlau,
% diese Farbdefinitionen sollten für den Druck verwendet werden (alles schwarz)
    %linkcolor=black, % einfache interne Verknüpfungen
    %anchorcolor=black, % Ankertext
    %citecolor=black, % Verweise auf Literaturverzeichniseinträge im Text
    %filecolor=black, % Verknüpfungen, die lokale Dateien öffnen
    %menucolor=black, % Acrobat-Menüpunkte
    %urlcolor=black,
%
    %backref, % Quellen werden zurück auf ihre Zitate verlinkt
    pdftex,
    plainpages=false, % zur korrekten Erstellung der Bookmarks
    pdfpagelabels=true, % zur korrekten Erstellung der Bookmarks
    hypertexnames=false, % zur korrekten Erstellung der Bookmarks
    linkcolor=black,
    linktoc=all,
]{hyperref}
% Befehle, die Umlaute ausgeben, führen zu Fehlern, wenn sie hyperref als Optionen übergeben werden
\hypersetup{
    pdftitle={\titel -- \untertitel},
    pdfauthor={\autorName},
    pdfcreator={\autorName},
    pdfsubject={\titel -- \untertitel},
    pdfkeywords={\titel -- \untertitel},
}


% zum Einbinden von Programmcode -----------------------------------------------
\usepackage{listings}
\usepackage{xcolor}
\definecolor{hellgelb}{rgb}{1,1,0.9}
\definecolor{colKeys}{rgb}{0,0,1}
\definecolor{colIdentifier}{rgb}{0,0,0}
\definecolor{colComments}{rgb}{0,0.5,0}
\definecolor{colString}{rgb}{1,0,0}
\lstset{
    float=hbp,
	basicstyle=\footnotesize,
    identifierstyle=\color{colIdentifier},
    keywordstyle=\color{colKeys},
    stringstyle=\color{colString},
    commentstyle=\color{colComments},
    backgroundcolor=\color{hellgelb},
    columns=flexible,
    tabsize=2,
    frame=single,
    extendedchars=true,
    showspaces=false,
    showstringspaces=false,
    numbers=left,
    numberstyle=\tiny,
    breaklines=true,
    breakautoindent=true,
	captionpos=b,
}
\lstdefinelanguage{cs}{
	sensitive=false,
	morecomment=[l]{//},
	morecomment=[s]{/*}{*/},
	morestring=[b]",
	morekeywords={
		abstract,event,new,struct,as,explicit,null,switch
		base,extern,object,this,bool,false,operator,throw,
		break,finally,out,true,byte,fixed,override,try,
		case,float,params,typeof,catch,for,private,uint,
		char,foreach,protected,ulong,checked,goto,public,unchecked,
		class,if,readonly,unsafe,const,implicit,ref,ushort,
		continue,in,return,using,decimal,int,sbyte,virtual,
		default,interface,sealed,volatile,delegate,internal,short,void,
		do,is,sizeof,while,double,lock,stackalloc,
		else,long,static,enum,namespace,string},
}
\lstdefinelanguage{natural}{
	sensitive=false,
	morecomment=[l]{/*},
	morestring=[b]",
	morestring=[b]',
	alsodigit={-,*},
	morekeywords={
		DEFINE,DATA,LOCAL,END-DEFINE,WRITE,CALLNAT,PARAMETER,USING,
		IF,NOT,END-IF,ON,*ERROR-NR,ERROR,END-ERROR,ESCAPE,ROUTINE,
		PERFORM,SUBROUTINE,END-SUBROUTINE,CONST,END-FOR,END,FOR,RESIZE,
		ARRAY,TO,BY,VALUE,RESET,COMPRESS,INTO,EQ},
}
\lstdefinelanguage{php}{
	sensitive=false,
	morecomment=[l]{/*},
	morestring=[b]",
	morestring=[b]',
	alsodigit={-,*},
	morekeywords={
		abstract,and,array,as,break,case,catch,cfunction,class,clone,const,
		continue,declare,default,do,else,elseif,enddeclare,endfor,endforeach,
		endif,endswitch,endwhile,extends,final,for,foreach,function,global,
		goto,if,implements,interface,instanceof,namespace,new,old_function,or,
		private,protected,public,static,switch,throw,try,use,var,while,xor
		die,echo,empty,exit,eval,include,include_once,isset,list,require,
		require_once,return,print,unset},
}
 % verwendete Packages
% !TEX root = ../Projektdokumentation.tex

% Seitenränder -----------------------------------------------------------------
\setlength{\topskip}{\ht\strutbox} % behebt Warnung von geometry
\geometry{a4paper,left=20mm,right=20mm,top=25mm,bottom=35mm}

\usepackage[
	automark, % Kapitelangaben in Kopfzeile automatisch erstellen
	headsepline, % Trennlinie unter Kopfzeile
	ilines % Trennlinie linksbündig ausrichten
]{scrlayer-scrpage}

% Kopf- und Fußzeilen ----------------------------------------------------------
\pagestyle{scrheadings}
% chapterpagestyle gibt es nicht in scrartcl
%\renewcommand{\chapterpagestyle}{scrheadings}
\clearscrheadfoot

% Kopfzeile
\renewcommand{\headfont}{\normalfont} % Schriftform der Kopfzeile
\ihead{\large{\textsc{\titel}}\\ \small{\untertitel} \\[2ex] \textit{\headmark}}
\chead{}
\ohead{\includegraphics[scale=0.2]{\betriebLogo}}
\setlength{\headheight}{15mm} % Höhe der Kopfzeile
%\setheadwidth[0pt]{textwithmarginpar} % Kopfzeile über den Text hinaus verbreitern (falls Logo den Text überdeckt)

% Fußzeile
\ifoot{\autorName}
\cfoot{}
\ofoot{\pagemark}

% Überschriften nach DIN 5008 in einer Fluchtlinie
% ------------------------------------------------------------------------------

% Abstand zwischen Nummerierung und Überschrift definieren
% > Schön wäre hier die dynamische Berechnung des Abstandes in Abhängigkeit
% > der Verschachtelungstiefe des Inhaltsverzeichnisses
\newcommand{\headingSpace}{1.5cm}

% Abschnittsüberschriften im selben Stil wie beim Inhaltsverzeichnis einrücken
\renewcommand*{\othersectionlevelsformat}[3]{
  \makebox[\headingSpace][l]{#3\autodot}
}

% Für die Einrückung wird das Paket tocloft benötigt
%\cftsetindents{chapter}{0.0cm}{\headingSpace}
\cftsetindents{section}{0.0cm}{\headingSpace}
\cftsetindents{subsection}{0.0cm}{\headingSpace}
\cftsetindents{subsubsection}{0.0cm}{\headingSpace}
\cftsetindents{figure}{0.0cm}{\headingSpace}
\cftsetindents{table}{0.0cm}{\headingSpace}


% Allgemeines
% ------------------------------------------------------------------------------

\onehalfspacing % Zeilenabstand 1,5 Zeilen
\frenchspacing % erzeugt ein wenig mehr Platz hinter einem Punkt

% Schusterjungen und Hurenkinder vermeiden
\clubpenalty = 10000
\widowpenalty = 10000
\displaywidowpenalty = 10000

% Quellcode-Ausgabe formatieren
\lstset{numbers=left, numberstyle=\tiny, numbersep=5pt, breaklines=true}
\lstset{emph={square}, emphstyle=\color{red}, emph={[2]root,base}, emphstyle={[2]\color{blue}}}

\counterwithout{footnote}{section} % Fußnoten fortlaufend durchnummerieren
\setcounter{tocdepth}{3} % im Inhaltsverzeichnis werden die Kapitel bis zum Level der subsubsection übernommen
\setcounter{secnumdepth}{3} % Kapitel bis zum Level der subsubsection werden nummeriert

% Aufzählungen anpassen
\renewcommand{\labelenumi}{\arabic{enumi}.}
\renewcommand{\labelenumii}{\arabic{enumi}.\arabic{enumii}.}
\renewcommand{\labelenumiii}{\arabic{enumi}.\arabic{enumii}.\arabic{enumiii}}

% Tabellenfärbung:
\definecolor{heading}{rgb}{0.64,0.78,0.86}
\definecolor{odd}{rgb}{0.9,0.9,0.9}
 % Definitionen zum Aussehen der Seiten
% !TEX root = ../Projektdokumentation.tex

% Abkürzungen, ggfs. mit korrektem Leerraum
\newcommand{\bs}{$\backslash$\xspace}
\newcommand{\bspw}{bspw.\xspace}
\newcommand{\bzw}{bzw.\xspace}
\newcommand{\ca}{ca.\xspace}
\newcommand{\dahe}{\mbox{d.\,h.}\xspace}
\newcommand{\etc}{etc.\xspace}
\newcommand{\eur}[1]{\mbox{#1\,\texteuro}\xspace}
\newcommand{\evtl}{evtl.\xspace}
\newcommand{\ggfs}{ggfs.\xspace}
\newcommand{\Ggfs}{Ggfs.\xspace}
\newcommand{\gqq}[1]{\glqq{}#1\grqq{}}
\newcommand{\inkl}{inkl.\xspace}
\newcommand{\insb}{insb.\xspace}
\newcommand{\ua}{\mbox{u.\,a.}\xspace}
\newcommand{\usw}{usw.\xspace}
\newcommand{\Vgl}{Vgl.\xspace}
\newcommand{\zB}{\mbox{z.\,B.}\xspace}

% Befehle für häufig anfallende Aufgaben
\newcommand{\Abbildung}[1]{\autoref{fig:#1}}
\newcommand{\Anhang}[1]{\appendixname{}~\ref{#1}: \nameref{#1} \vpageref{#1}}
\newcommand{\includegraphicsKeepAspectRatio}[2]{\includegraphics[width=#2\textwidth,height=#2\textheight,keepaspectratio]{#1}}
\newcommand{\Zitat}[2][\empty]{\ifthenelse{\equal{#1}{\empty}}{\citep{#2}}{\citep[#1]{#2}}}
\newcommand{\Autor}[1]{\textsc{#1}} % zum Ausgeben von Autoren
\newcommand{\itemd}[2]{\item{\textbf{#1}}\\{#2}} % erzeugt ein Listenelement mit fetter Überschrift

% fügt Tabellen aus einer TEX-Datei ein
\newcommand{\tabelle}[3] % Parameter: caption, label, file
{\begin{table}[htbp]
\centering
\singlespacing
\input{Tabellen/#3}
\caption{#1}
\label{#2}
\end{table}}

\newcommand{\tabelleAnhang}[1] % Parameter: file
{\begin{center}
\singlespacing
\input{Tabellen/#1}
\end{center}}

% einfaches Wechseln der Schrift, z.B.: \changefont{cmss}{sbc}{n}
\newcommand{\changefont}[3]{\fontfamily{#1} \fontseries{#2} \fontshape{#3} \selectfont}

% Verwendung analog zu \includegraphics
\newlength{\myx} % Variable zum Speichern der Bildbreite
\newlength{\myy} % Variable zum Speichern der Bildhöhe
\newcommand\includegraphicstotab[2][\relax]{%
% Abspeichern der Bildabmessungen
\settowidth{\myx}{\includegraphics[{#1}]{#2}}%
\settoheight{\myy}{\includegraphics[{#1}]{#2}}%
% das eigentliche Einfügen
\parbox[c][1.1\myy][c]{\myx}{%
\includegraphics[{#1}]{#2}}%
}

\definecolor{AOBlau}{rgb}{0, 0.28, 0.56}

% verschiedene Befehle um Wörter semantisch auszuzeichnen ----------------------
\newcommand{\Index}[2][\empty]{\ifthenelse{\equal{#1}{\empty}}{\index{#2}#2}{\index{#1}#2}}
\newcommand{\Fachbegriff}[2][\empty]{\ifthenelse{\equal{#1}{\empty}}{\textit{\Index{#2}}}{\textit{\Index[#1]{#2}}}}
\newcommand{\NeuerBegriff}[2][\empty]{\ifthenelse{\equal{#1}{\empty}}{\textbf{\Index{#2}}}{\textbf{\Index[#1]{#2}}}}

\newcommand{\Ausgabe}[1]{\texttt{#1}}
\newcommand{\Eingabe}[1]{\texttt{#1}}
\newcommand{\Code}[1]{\texttt{#1}}
\newcommand{\Datei}[1]{\texttt{#1}}

\newcommand{\Assembly}[1]{\textsf{#1}}
\newcommand{\Klasse}[1]{\textsf{#1}}
\newcommand{\Methode}[1]{\textsf{#1}}
\newcommand{\Attribut}[1]{\textsf{#1}}

\newcommand{\Datentyp}[1]{\textsf{#1}}
\newcommand{\XMLElement}[1]{\textsf{#1}}
\newcommand{\Webservice}[1]{\textsf{#1}}

\newcommand{\Refactoring}[1]{\Fachbegriff{#1}}
\newcommand{\CodeSmell}[1]{\Fachbegriff{#1}}
\newcommand{\Metrik}[1]{\Fachbegriff{#1}}
\newcommand{\DesignPattern}[1]{\Fachbegriff{#1}}
 % eigene allgemeine Befehle, die z.B. die Arbeit mit LaTeX erleichtern
% Abkürzungen
\newcommand{\Versis}{\textsc{Versis}\xspace}
\newcommand{\NI}{NatInfo\xspace}
\newcommand{\AO}{\textsc{Alte Oldenburger} Krankenversicherung\xspace}
 % eigene projektspezifische Befehle, z.B. Abkürzungen usw.

\begin{document}

\phantomsection
\thispagestyle{plain}
\pdfbookmark[1]{Deckblatt}{deckblatt}
% !TEX root = Projektdokumentation.tex
\begin{titlepage}

\begin{center}
\includegraphics[scale=0.08]{ihk-dortmund.jpeg}\\[1ex]
\Large{Abschlussprüfung \pruefungstermin}\\[3ex]

\Large{\ausbildungsberuf}\\
\LARGE{\betreff}\\[4ex]

\huge{\textbf{\titel}}\\[1.5ex]
\Large{\textbf{\untertitel}}\\[4ex]

\normalsize
Abgabetermin: \abgabeOrt, den \abgabeTermin\\[3em]
\textbf{Prüfungsbewerber:}\\
\autorName\\
\autorAnschrift\\
\autorOrt\\[5ex]

\includegraphics[scale=1]{\betriebLogo}\\[2ex]
\textbf{Ausbildungsbetrieb:}\\
\betriebName\\
\betriebAnschrift\\
\betriebOrt\\[5em]
\end{center}

\small
\noindent
Dieses Werk einschließlich seiner Teile ist \textbf{urheberrechtlich geschützt}.
Jede Verwertung außerhalb der engen Grenzen des Urheberrechtgesetzes ist ohne
Zustimmung des Autors unzulässig und strafbar. Das gilt insbesondere für
Vervielfältigungen, Übersetzungen, Mikroverfilmungen sowie die Einspeicherung
und Verarbeitung in elektronischen Systemen.

\end{titlepage}
\cleardoublepage

% Preface --------------------------------------------------------------------
\phantomsection
\pagenumbering{Roman}
\pdfbookmark[1]{Inhaltsverzeichnis}{inhalt}
\tableofcontents

\cleardoublepage

\phantomsection
\listoffigures
\cleardoublepage

\phantomsection
\listoftables
\cleardoublepage

\newcommand{\abkvz}{Abkürzungsverzeichnis}
\renewcommand{\nomname}{\abkvz}
\section*{\abkvz}
\markboth{\abkvz}{\abkvz}
\addcontentsline{toc}{section}{\abkvz}
% !TEX root = Projektdokumentation.tex

% Es werden nur die Abkürzungen aufgelistet, die mit \ac definiert und auch benutzt wurden. 
%
% \acro{VERSIS}{Versicherungsinformationssystem\acroextra{ (Bestandsführungssystem)}}
% Ergibt in der Liste: VERSIS Versicherungsinformationssystem (Bestandsführungssystem)
% Im Text aber: \ac{VERSIS} -> Versicherungsinformationssystem (VERSIS)

% Hinweis: allgemein bekannte Abkürzungen wie z.B. bzw. u.a. müssen nicht ins Abkürzungsverzeichnis aufgenommen werden
% Hinweis: allgemein bekannte IT-Begriffe wie Datenbank oder Programmiersprache müssen nicht erläutert werden,
%          aber ggfs. Fachbegriffe aus der Domäne des Prüflings (z.B. Versicherung)

% Die Option (in den eckigen Klammern) enthält das längste Label oder
% einen Platzhalter der die Breite der linken Spalte bestimmt.
\begin{acronym}[WWWWW]
	\acro{AO}{\textsc{Alte Oldenburger} Krankenversicherung AG}
	\acro{API}{Application Programming Interface}
	\acro{CSS}{Cascading Style Sheets}
	\acro{CSV}{Comma Separated Value}
	\acro{EPK}{Ereignisgesteuerte Prozesskette}
	\acro{ERM}{En\-ti\-ty-Re\-la\-tion\-ship-Mo\-dell}
	\acro{HTML}{Hypertext Markup Language}\acused{HTML}
	\acro{IDE}{Integrated Development Environment}
	\acro{MVC}[MVC]{Model View Controller}
	\acro{NatInfo}[\textsc{NatInfo}]{Natural Information System}
	\acro{Natural}[\textsc{Natural}]{Programmiersprache der Software AG}\acused{Natural}
	\acro{ORM}{Object-Relational Mapping}
	\acro{PHP}{Hypertext Preprocessor}
	\acro{SDK}{Software Development Kit}
	\acro{SQL}{Structured Query Language}
	\acro{SVN}{Subversion}
	\acro{UML}{Unified Modeling Language}
	\acro{XML}{Extensible Markup Language}
\end{acronym}

\clearpage

% Inhalt ---------------------------------------------------------------------
\pagenumbering{arabic}
% !TEX root = Projektdokumentation.tex
% !TEX root = ../Projektdokumentation.tex
\section{Einleitung}
\label{sec:Einleitung}


\subsection{Projektumfeld} 
\label{sec:Projektumfeld}
Die bits \& likes GmbH(B \& L) ist eine Fullservice-Digitalagentur aus Dortmund. Sie bietet sowohl Onlinemarketing als auch Websitenentwicklung an. Aktuell beschäfftigt sie ca. 50 Mitarbeiter, die Tendenz ist steigend.

Die Idee für das Projekt kam bei der Entwicklung einer Website für die Murtfeldt Kunststoffe GmbH \& Co. KG(Murtfeldt). Historisch nutzte Murtfeldt Typo3 als Content-Management-System. B \& L entwickelte mit Nuxt2 eine Headless Lösung für diese neue Website. Dieser Technologiestack soll nun mit der neuesten Version von Typo3(Version 11.5 zum Zeitpunkt des Projekts) als auch der neusten Version von Nuxt neu entwickelt werden. Damit zukünftige Murtfeldt Projekt also auch andere Projekte(z.B. B \& L Firmenwebsite) eine Vorlage haben, mit welcher sie umgesetzt werden können.


\subsection{Projektziel} 
\label{sec:Projektziel}
Das Ziel des Projektes ist es eine Typo3 Nuxt3 Vorlage(Skeleton) zu erstellen, auf dessen Basis neue Projekte umgesetzt werden können. Dazu müssen folgende Punkte erreicht werden:
\begin{itemize}
	\item Content aus dem Typo3 System im JSON-Format ausspielen.
	\item Potentielle Erweiterungen an den JSON-Daten um alle technischen Voraussetzungen zu erfüllen
	\item Content im Nuxt3 Frontend auslesen, verarbeiten und auspielen. Dies beinhaltet sowohl den Content, als auch die Navigation
\end{itemize}

\subsection{Projektbegründung} 
\label{sec:Projektbegruendung}
Bis jetzt klonte B \& L immer alte Websiten auf dem Typo3 \& Nuxt2 Technologiestack um neue Websiten zu entwickeln. Dies bedeutete, dass viel Content, als auch alter Code erstmal händisch aus dem Projekt gelöscht werden musste. Zusätzlich musste die Datenbank bereinigt werden, damit die neue Website, keine Altlasten / kritische Daten aus dem vorherigen Projekt mit übernimmt. Die Alternative war, ein komplett neues Typo3 und Nuxt2 Projekt aufzusetzen und jedes mal viel Code / Template-Anpassungen neu schreiben. Dieser unnötige Aufwand soll mit diesem Projekt umgangen werden. Dadurch entstehen sowohl Kostenersparnisse durch kürzere Entwicklungszeiten, als auch weniger potentielle Fehler/Komplikationen mit alten Code/Daten von anderen Projekten.

\subsection{Projektschnittstellen} 
\label{sec:Projektschnittstellen}

Die finale Nuxt Anwendung interagiert nur mit der Typo3 Schnittstelle. Projekte, die mit dem Skeleton entwickelt werden, können aber weitere Schnittstellen einbinden. Im Fall von B \& L würde dies beispielsweise oft eine Schnittstelle zu einem Shopware-System(E-Commerce Platform) sein. Zukünftige Iterationen des Skeletons könnten diese auch als Standard beinhalten.

Die Mittel für das Projekt werden von  B \& L zur Verfügung gestellt. Da die Entwickler von B \& L auch die zukünftigen Nutzer des Projekts sind. Dementsprechend muss das Projekt auch den Entwicklern später präsentiert und eine Dokumentation für diese geschrieben werden. Projekte, welche aus dem Skeleton entwickelt werden, würden aber in der Zukunft Kunden präsentiert werden und Teile des Skeletons enthalten.


\subsection{Projektabgrenzung} 
\label{sec:Projektabgrenzung}

Das Projekte enthält folgende features nicht:

\begin{itemize}
	\item Ein umgesetztes Frontend-Design
	\item Alle gängigen Typo3 Komponenten im Frontend eingebaut(nur die Basics)
	\item Anbindungen zu anderen technischen Schnittstellen außer Typo3
\end{itemize}

% !TEX root = ../Projektdokumentation.tex
\section{Projektplanung} 
\label{sec:Projektplanung}


\subsection{Projektphasen}
\label{sec:Projektphasen}

Die Bearbeitung der Projektarbeit inklusive Dokumentation fand vom 23.09.2022 - 11.12.2022 statt. Die Planungs- und Implementierungsphasen schließen im Oktober ab. Die Dokumentation wurde in den folgenden Monaten abgeschlossen.

Die im Projektantrag geplante Zeitplanung sah wie folgt aus:

Tabelle~\ref{tab:Zeitplanung} zeigt die alte grobe Zeitplanung des Projektes.
\tabelle{alte grobe Zeitplanung}{tab:Zeitplanung}{ZeitplanungKurz}


Es wurden aber Änderungen an der Zeitplanung vorgenommen, die neue Zeitplanung sieht wie folgt aus:

Tabelle~\ref{tab:Zeitplanung} zeigt die neue grobe Zeitplanung des Projektes.
\tabelle{neue grobe Zeitplanung}{tab:Zeitplanung}{ZeitplanungKurzNeu}

Warum vom Projektantrag abgewichen wurden, wird in Kapitel \ref{sec:AbweichungenProjektantrag} Abweichungen vom Projektantrag erklärt.

\subsubsection{Implementierungsphase Zeitplanung}
\label{sec:Implementierungsphase Zeitplanung}

Die Implementierungsphase nimmt den mit Abstand größten Teil der Projektarbeit ein. Dadurch ist es sinnvoll hier eine detailierte Zeitplanung zu erstellen.

Tabelle~\ref{tab:Zeitplanung Implementierungsphase} zeigt die Zeitplanung der Implementierungsphase.
\tabelle{Zeitplanung Implementierungsphase}{tab:Zeitplanung Implementierungsphase}{ZeitplanungImplementierungsphase}

Wie zu sehen ist, wird geplant, dass die Nuxt 3 Entwicklungen 80\% des Arbeitsumfangs ausmachen werden. Dies liegt daran, dass ein großer Teil der Arbeitsaufwände durch das installieren der Headless Extension bereits erledigt werden. Der Prüfungsbewerber hofft zusätzlich, dass er die Erweiterungen des Typo3 teilweise aus alten Projekten übernehmen kann. Ob diese mit der neusten Typo3 Version funktionieren, wird sich im Verlaufe des Projekts zeigen. Falls Code aus alten Projekten recycled wird, wird dies im Text erkennbar gemacht. Dementsprechend wird sich Kapitel \ref{sec:Implementierungsphase} Implementierungsphase auch primär auf die Nuxt 3 Entwicklungen fokussieren.

\subsection{Abweichungen vom Projektantrag}
\label{sec:AbweichungenProjektantrag}

An dem Projekt als solchem hat sich nichts geändert. Es wurde das gleiche Stück Software erstellt, wie es im Projektantrag beschrieben wurde. Externe Umstände und fehlerhafte Planung haben es leider notwendig gemacht, dass Anpassungen an dem Bearbeitungszeitraum und der Zeitplanung getroffen werden mussten.

\subsubsection{Bearbeitungszeitraum}
\label{sec:AbweichungenProjektantrag}
Der geplante Bearbeitungszeitraum wurde relativ schnell nach hinten verschoben. Der Grund waren mehrere Covid-19 Erkrankungen im Betrieb. Der Prüfungsbewerber selber wurde ebenfalls während des geplanten Zeitraums für ein paar Tage krank. Dadurch konnte der Prüfungsbewerber weniger Zeit in die Projektarbeit stecken als geplant, weswegen das Ende der Bearbeitungszeit in den Dezember verschoben wurde. Ursprünglich war geplant, dass der Bearbeitungszeitraum am 20.09.2022 startet und am 15.11.2022 endet. Der ursprüngliche Bearbeitungszeitraum wurde glücklicherweise so bewusst gewählt, dass er bei Notfällen etwas verlängert werden kann. Der finale Abgabetermin war der 14.12.2022 und konnte eingehalten werden.

\subsubsection{Zeitplanung}
\label{sec:AbweichungenProjektantrag}

Kurz nach einreichen des Projektantrages, wurde dem Prüfungsbewerber relativ schnell klar, dass die Zeitplanung so nicht funktionieren würde. Er hatte im Projektantrag nur 5 Stunden für die Dokumentation des Projektes eingeplant. Das entspricht weniger als 9 Minuten pro Seite(300 Minuten / 35 Seiten). In so einem kurzem Zeitraum lässt sich keine angemessene Dokumentation erstellen. Der Prüfungsbewerber ist ursprünglich davon ausgegangen, dass der Projektantrag, aufgrund dieser Zeit, abgelehnt werden würde. Dies ist nicht geschehen. Da er aber unmöglich diese Dokumentation in 5 Stunden schreiben konnte, wurden die Zeiten der Projektphasen abgeändert. Dafür wurden 15 Stunden aus der Implementierungsphase entnommen. Diese Phase hat mit großem Abstand die meisten Stunden. Selbst mit Entnahme der Stunden, hatte diese Phase immer noch die Hälfte der gesamten Projektzeit. Hier wurden bei der Planung des Projekts die Zeiten auch am großzügigsten geschätzt. 5 weitere Stunden wurden aus der Analysephase entnommen. Da ein guter Teil der Arbeit in dieser Phase bereits beim einreichen des Projektantrages erledigt wurde. So hatte der Prüfungsbewerber nun 25 Stunden Zeit die Dokumentation zu schreiben.

Ob die neue Planung funktioniert hat, oder ob es während des Projektes noch zu anderen Problemen kam, wird im Fazit analysiert und besprochen.

\subsection{Ressourcenplanung}
\label{sec:Ressourcenplanung}

Für die Umsetzung des Projektes wurden folgende Ressourcen benötigt:
\begin{itemize}
	\item Computer
	\item Internetzugang
	\item Büroraum
	\item Quellcode von Typo3 und Nuxt 3 (kostenlos übers Internet verfügbar)
	\item Personelle Ressourcen: senior Developer für Rückfragen
\end{itemize}

Zusätzlich wurde eine LaTeX-Vorlage zum erstellen der Dokumentation verwendet.  Die LaTeX-Vorlage\footnote{\Vgl https://it-berufe-podcast.de/vorlage-fuer-die-projektdokumentation/} wurde von Stefan Macke\footnote{\Vgl Blog des Autors: http://fachinformatiker-anwendungsentwicklung.net, Twitter: @StefanMacke.} entwickelt. Die Vorlage steht unter der Creative Commons Namensnennung Lizenz. Die Vorlage selbst hat keinen Inhalt. Sie dient als eine grobe Strukturierung der Dokumentation.

\subsection{Entwicklungsprozess}
\label{sec:Entwicklungsprozess}

Bei der Entwicklung des Projektes wurde ein Wasserfall Entwicklungsprozess benutzt. Dies bedeutet dass die Projektphasen linear und ohne Rückschritte nacheinander abgearbeitet wurden. Die vordefinierten Projektphasen ergaben sich aus dem Projektantrag.

% !TEX root = ../Projektdokumentation.tex
\section{Analysephase} 
\label{sec:Analysephase}


\subsection{Ist-Analyse} 
\label{sec:IstAnalyse}
\begin{itemize}
	\item Wie ist die bisherige Situation (\zB bestehende Programme, Wünsche der Mitarbeiter)?
	\item Was gilt es zu erstellen/verbessern?
\end{itemize}


\subsection{Wirtschaftlichkeitsanalyse}
\label{sec:Wirtschaftlichkeitsanalyse}
\begin{itemize}
	\item Lohnt sich das Projekt für das Unternehmen?
\end{itemize}


\subsubsection{\gqq{Make or Buy}-Entscheidung}
\label{sec:MakeOrBuyEntscheidung}
\begin{itemize}
	\item Gibt es vielleicht schon ein fertiges Produkt, dass alle Anforderungen des Projekts abdeckt?
	\item Wenn ja, wieso wird das Projekt trotzdem umgesetzt?
\end{itemize}


\subsubsection{Projektkosten}
\label{sec:Projektkosten}
\begin{itemize}
	\item Welche Kosten fallen bei der Umsetzung des Projekts im Detail an (\zB Entwicklung, Einführung/Schulung, Wartung)?
\end{itemize}

\paragraph{Beispielrechnung (verkürzt)}
Die Kosten für die Durchführung des Projekts setzen sich sowohl aus Personal-, als auch aus Ressourcenkosten zusammen.
Laut Tarifvertrag verdient ein Auszubildender im dritten Lehrjahr pro Monat \eur{1000} Brutto. 

\begin{eqnarray}
8 \mbox{ h/Tag} \cdot 220 \mbox{ Tage/Jahr} = 1760 \mbox{ h/Jahr}\\
\eur{1000}\mbox{/Monat} \cdot 13,3 \mbox{ Monate/Jahr} = \eur{13300} \mbox{/Jahr}\\
\frac{\eur{13300} \mbox{/Jahr}}{1760 \mbox{ h/Jahr}} \approx \eur{7,56}\mbox{/h}
\end{eqnarray}

Es ergibt sich also ein Stundenlohn von \eur{7,56}. 
Die Durchführungszeit des Projekts beträgt 70 Stunden. Für die Nutzung von Ressourcen\footnote{Räumlichkeiten, Arbeitsplatzrechner etc.} wird 
ein pauschaler Stundensatz von \eur{15} angenommen. Für die anderen Mitarbeiter wird pauschal ein Stundenlohn von \eur{25} angenommen. 
Eine Aufstellung der Kosten befindet sich in Tabelle~\ref{tab:Kostenaufstellung} und sie betragen insgesamt \eur{2739,20}.
\tabelle{Kostenaufstellung}{tab:Kostenaufstellung}{Kostenaufstellung.tex}


\subsubsection{Amortisationsdauer}
\label{sec:Amortisationsdauer}
\begin{itemize}
	\item Welche monetären Vorteile bietet das Projekt (\zB Einsparung von Lizenzkosten, Arbeitszeitersparnis, bessere Usability, Korrektheit)?
	\item Wann hat sich das Projekt amortisiert?
\end{itemize}

\paragraph{Beispielrechnung (verkürzt)}
Bei einer Zeiteinsparung von 10 Minuten am Tag für jeden der 25 Anwender und 220 Arbeitstagen im Jahr ergibt sich eine gesamte Zeiteinsparung von 
\begin{eqnarray}
25 \cdot 220 \mbox{ Tage/Jahr} \cdot 10 \mbox{ min/Tag} = 55000 \mbox{ min/Jahr} \approx 917 \mbox{ h/Jahr} 
\end{eqnarray}

Dadurch ergibt sich eine jährliche Einsparung von 
\begin{eqnarray}
917 \mbox{h} \cdot \eur{(25 + 15)}{\mbox{/h}} = \eur{36680}
\end{eqnarray}

Die Amortisationszeit beträgt also $\frac{\eur{2739,20}}{\eur{36680}\mbox{/Jahr}} \approx 0,07 \mbox{ Jahre} \approx 4 \mbox{ Wochen}$.


\subsection{Nutzwertanalyse}
\label{sec:Nutzwertanalyse}
\begin{itemize}
	\item Darstellung des nicht-monetären Nutzens (\zB Vorher-/Nachher-Vergleich anhand eines Wirtschaftlichkeitskoeffizienten). 
\end{itemize}

\paragraph{Beispiel}
Ein Beispiel für eine Entscheidungsmatrix findet sich in Kapitel~\ref{sec:Architekturdesign}: \nameref{sec:Architekturdesign}.


\subsection{Anwendungsfälle}
\label{sec:Anwendungsfaelle}
\begin{itemize}
	\item Welche Anwendungsfälle soll das Projekt abdecken?
	\item Einer oder mehrere interessante (!) Anwendungsfälle könnten exemplarisch durch ein Aktivitätsdiagramm oder eine \ac{EPK} detailliert beschrieben werden. 
\end{itemize}

\paragraph{Beispiel}
Ein Beispiel für ein Use Case-Diagramm findet sich im \Anhang{app:UseCase}.


\subsection{Qualitätsanforderungen}
\label{sec:Qualitaetsanforderungen}
\begin{itemize}
	\item Welche Qualitätsanforderungen werden an die Anwendung gestellt (\zB hinsichtlich Performance, Usability, Effizienz \etc (siehe \citet{ISO9126}))?
\end{itemize}


\subsection{Lastenheft/Fachkonzept}
\label{sec:Lastenheft}
\begin{itemize}
	\item Auszüge aus dem Lastenheft/Fachkonzept, wenn es im Rahmen des Projekts erstellt wurde.
	\item Mögliche Inhalte: Funktionen des Programms (Muss/Soll/Wunsch), User Stories, Benutzerrollen
\end{itemize}

\paragraph{Beispiel}
Ein Beispiel für ein Lastenheft findet sich im \Anhang{app:Lastenheft}. 

% !TEX root = ../Projektdokumentation.tex
\section{Entwurfsphase} 
\label{sec:Entwurfsphase}

\subsection{Zielplattform}
\label{sec:Zielplattform}

\subsubsection{TypeScript}
\label{sec:TypeScript}

Als Programmiersprache wurde TypeScript(\acs{TS}) ausgewählt\footnote{\Vgl \citet{TypeScript}.}. \acs{TS} ist ein Superset von JavaScript(\acs{JS}). Dies bedeutet, dass jeder JS-Code in \acs{TS} funktioniert, aber nicht jeder \acs{TS}-Code in JS funktioniert. \acs{TS} erweitert JS um Typen. Dies bedeutet, dass Variablen feste Datentypen(string, number, etc.) haben können. Dies hat den Vorteil, dass \acs{TS} beim kompilieren Fehler schmeißt, falls eine Variable einen Wert bekommt, welche sie nicht haben sollte. \acs{TS} wird zu JS kompiliert, da Browser aktuell nur JS und WebAssembly unterstützen. Dadurch werden vor allem große Projekte weniger fehleranfällig. Das Skeleton selber ist nicht so ein großes Projekt, aber es werden große Projekte damit entwickelt werden. Deswegen ist es sinnvoll, bereits mit \acs{TS} zu starten.

\subsubsection{Nuxt 3}
\label{sec:Nuxt 3}

Als Framework für TS wurde sich für Nuxt 3 entschieden. Nuxt 3 ist ein serverseitig gerendertes Javascript-Framework mit TypeScript Support. Nuxt 3 bietet viele Features, welche Webentwicklung einfacher machen. Dazu zählt beispielsweise reactivty. Dies bedeutet, dass das Framework automatisch Variablen in der Darstellung anpassen kann, wenn diese ihren Wert ändern. Viele andere Javascript-Frameworks werden clientseitig gerendert. Das bedeutet, dass der Nutzer das HTML erst rendern muss. Nuxt3 macht dies bereits auf dem Server, wodurch die Seite beim Nutzer schneller angezeigt wird. Dies wirkt sich auch vorteilhaft auf SEO aus.

\subsubsection{DaisyUI / Tailwind}
\label{sec:DaisyUI / Tailwind}

Als UI-Framework wurde DaisyUI genutzt. Für die Umsetzung des Skeletons war die Nutzung eines UI-Frameworks nicht wirklich nötig. Insgesamt musste während der Entwicklung des Skeleton nicht viel CSS / Styling genutzt werden. Das Framework liefert vorgestylte Komponenten und CSS-Klassen, welche die zukünftige Entwicklung beschleunigen sollen. DaisyUI basiert auf Tailwind, weswegen alle Tailwind Klassen auch in dem Projekt genutzt werden können.

\subsubsection{Docker}
\label{sec:Docker}

Um Typo3 besser testen zu können, wird Typo3 in einem Docker-Container ausgeführt. Docker ist eine Open-Source-Tool, welches es erlaubt, Software in abgekapselten Umgebungen/Containern laufen zu lassen. Ein Docker-Container ist Softwarepaket, was alle nötigen Werkzeuge enthält, um ein Programm laufen zu lassen. Docker ermöglicht es, verschiedenste Programme (andere Programmiersprachen, Systemanforderungen, Libraries, etc.) in unterschiedlichen Containern laufen zu lassen. Beispielsweise kann auf einem Rechner/Server in einem Docker-Container Typo3 laufen, während in einem anderen Nuxt läuft. Docker-Container können schnell hoch und runtergefahren werden, was testen von neuem Code, oder neustarten von Systemen sehr schnell macht. 

Bei Docker-Containern handelt es sich nicht um virtuelle Maschinen. Vielmehr sind Docker-Container virtuelle Betriebssysteme. Dadurch können Docker-Container mehr Ressourcen mit ihrem Host-System teilen. Deswegen verbrauchen sie insgesamt weniger Leistung als virtuelle Maschinen.

\subsubsection{Sonstiges}
\label{sec:Sonstiges}

Sonst wurden die klassischen Webtechnologien HTML und CSS genutzt.

Genauere Vergleiche zu anderen Frameworks / Technologien, werden im folgenden Kapitel analysiert.


\subsection{Architekturdesign}
\label{sec:Architekturdesign}

Anhand folgender Bewertungsmatrix, wurde sich für Nuxt 3 als JS-Framework entschieden:

\tabelle{Entscheidungsmatrix Nuxt 3}{tab:Entscheidungsmatrix Nuxt 3}{Nutzwert_Nuxt3.tex}

Als Optionen wurden alle Frameworks genommen, mit welchen \acs{BAL} bereits Erfahrungen gesammelt hat. Eigenentwicklungen wurden von Anfang an ausgeschlossen, da diese zu aufwendig sind. Andere Frameworks hätten eine zu lange Einarbeitung/Umschulung für potenziell zu wenig Nutzen gehabt.

Anhand folgender Bewertungsmatrix, wurde sich für DaisyUI als UI-Framework entschieden:

\tabelle{Entscheidungsmatrix DaisyUI}{tab:Entscheidungsmatrix DaisyUI}{Nutzwert_DaisyUI.tex}

Anhand folgender Bewertungsmatrix, wurde sich für TS als Programmiersprache entschieden:

\tabelle{Entscheidungsmatrix TS}{tab:Entscheidungsmatrix TS}{Nutzwert_TS.tex}

Beide Programmiersprachen werden von Nuxt 3 standardmäßg unterstützt. \acs{TS} benötigt mehr Zeilen Code und fügt einen weiteren Layer an Komplexität zu dem Projekt hinzu. Dafür werden Fehler früher entdeckt und das Projekt lässt sich leichter skalieren. Diese Faktoren werden von \acs{BAL} mehr wertgeschätzt, weswegen sich für \acs{TS} entschieden wurde.

Für das Typo3 musste keine Entscheidungsmatrix erstellt werden, da Typo3 von wichtigen Kunden vorgegeben wird. Es wird unabhängig vom Nutzen der Technologie ein Skeleton mit dieser Technologie benötigt. Typo3 hat unabhängig davon aber einige Vorteile, welche bereits erläutert wurden.

Navigation: Standardmäßig gibt die Headless-Extension von Typo3 keine Navigationsstruktur aus, weswegen noch Anpassungen am Template gemacht werden mussten. Es wurde sich entschieden, dass alle Seiten(und deren Unterseiten) in die Navigation ausgespielt werden, welche einer bestimmten Kategorie zugewiesen werden. Kategorien sind standardmäßig bei Seiten in Typo3 vorhanden und können mit kleinen Template-Anpassungen ausgespielt werden. Die Headless Extension bietet dafür Funktionalitäten.

Layout / Spalteneinteilung: Standardmäßig hat Typo3 keine Tools, um Komponenten in Spalten/Bereiche aufzuteilen. Es gibt viele Standardkonfigurationen, welche dies ermöglichen, aber ein komplett leeres Typo3 hat diese Funktionalität nicht. Deswegen wird eine eigene kleine Extension für Typo3 geschrieben, welche diese Funktionalität liefert.


\subsection{Entwurf der Benutzeroberfläche}
\label{sec:Benutzeroberflaeche} 

\paragraph{Frontend} 
\linebreak
Das Ziel der Benutzeroberfläche des Skeletons ist, dass Komponenten dargestellt werden und zwischen Seiten des Typo3 navigiert werden kann. Es muss noch keine gute Benutzeroberfläche erstellt werden, welche Kunden nützen können. Stattdessen muss es für Entwickler einfach und flexibel sein, eine gute Benutzeroberfläche für Kunden aus der Benutzeroberfläche des Skeletons, zu entwickeln. Dafür werden die Komponenten und die Navigation stilistisch sehr minimalistisch gehalten und primär ihre Logik erstellt. Die Navigation befindet sich klassisch am Anfang der Website und unter ihr wird der Content der Seite gerendert. Dies wurde auch hier angewandt.


\paragraph{Backend} 
\linebreak 
Die Benutzeroberfläche des Backends/CMS ist größtenteils durch Typo3 vorgegeben. Es werden lediglich die Komponenten erweitert, sodass gesteuert werden kann, wie viel Platz die Komponenten einnehmen. Dafür werden für die einzelnen Bildschirmgrößen(Desktop, Tablet, Mobile) Selects eingebaut, welche die Größen als Optionen haben.

\paragraph{Menüführung}  
\linebreak
Die Hauptnavigationspunkte, welche aus dem Typo3 ausgespielt werden, werden oben in der Menüleiste angezeigt. Ihre Kinder(Children/Unterseiten) werden ebenfalls ausgespielt, aber erst angezeigt, wenn durch die Navigation navigiert wird. Wenn einer der Hauptnavigationspunkte angeklickt wird, wird überprüft, ob dieser Children(Unterseiten) hat oder nicht. Wenn dies nicht der Fall ist, wird direkt die Seite von dem Hauptnavigationspunkt aufgerufen. Wenn dies der Fall ist, wird das Slidemenu aufgerufen. Das Slidemenu enthält sowohl den zuletzt angeklickten Navigationspunkt als auch alle Children von ihm. Es ist dazu da, um durch tiefere Navigationsebenen der Website zu steuern. Der Name des Slidemenu kommt von einer CSS Animation, welche das Menü in den Screen sliden lässt. 

Eine Darstellung der Menüführung findet sich im \Anhang{Screenshots des Menüs}. 

Zusätzlich gibt es einen Button, mit welchem zurück navigiert werden kann. Wenn der zuletzt angeklickte Navigationspunkt hier noch mal angeklickt wird, wird dessen Seite aufgerufen. Wenn ein Child angeklickt wird, wird wieder überprüft, ob es Children hat oder nicht. Falls nein, wird die Seite aufgerufen, falls ja wird das Child zum zuletzt angeklickten Navigationspunkt. 

\paragraph{rendern der verschiedenen Typo3-Inhaltselemente}
Im Typo3 können viele verschiedene Inhaltselemente gepflegt werden. Um diese darzustellen, wurde sich entschieden, für jedes Inhaltselement eine eigene Komponente zu erstellen. Für die Zuweisung der Inhaltselemente zu ihrer jeweiligen Nuxt 3 Komponente wurde eine weitere Nuxt 3 Komponente(TypoView.vue) erstellt. Diese analysiert/verarbeitet den Content der Seite und kümmert sich um die Zuweisungen.

\subsubsection{Typo3 Datenverarbeitung durch Nuxt / Flussdiagramm}
\label{sec:Typo3 Datenverarbeitung durch Nuxt / Flussdiagramm} 

Um die Typo3 Datenverarbeitung von Nuxt 3 darzustellen, wurde ein Flussdiagramm erstellt. siehe Abbildung \ref{FlussDiagramm} des Anhangs. Wie sich im Flussdiagramm erkennen lässt, werden alle Daten als erstes von einer Komponente([...slug].vue) ausgelesen. Diese Datei wertet teilweise die Daten aus, verschickt aber auch Teile der Daten weiter an andere Komponenten. 

Die Navigationsdaten im page Objekt werden an die Navigation.vue Komponente weitergeleitet. Die Navigation.vue Komponente besteht aus zwei weiteren Komponenten. Einmal der TopNavigation.vue und einmal der SlideMenu.vue. Alle drei Komponenten nutzen das navigation.ts Composable um miteinander zu kommunizieren und die Menüführung darzustellen.

Der Content wird an die Typoview.vue Komponente weitergeleitet. Dies ist die Komponente, welche sich um die Zuweisung der Typo3-Inhaltselemente zu den jeweiligen Nuxt 3 Komponenten kümmert. 

Die Breadcrumbs werden in die Breadcrumbs.vue Komponente weitergeleitet. Dort werden diese auch direkt dargestellt. 

Die Metadaten werden direkt in der [...slug].vue Komponente zu den Metainformationen der Website hinzugefügt.

\subsection{Maßnahmen zur Qualitätssicherung}
\label{sec:Qualitaetssicherung}

Um die SEO von dem Skeleton abschätzen zu können, kann Lighthouse von Google verwendet werden. Lighthouse analysiert die Seite bezüglich SEO und liefert abhängig davon einen Wert von 0 – 100. Falls es SEO Schwierigkeiten gibt, werden diese hier erläutert. Da noch nicht viel Content auf den Seiten gepflegt ist, kann es sein, dass es dadurch zu Abzügen kommt. Dies wird sich aber dann automatisch anpassen, wenn mit dem Skeleton eine vernünftige Website erstellt wird.

Wie gut der Code bezüglich Flexibilität und Leserlichkeit ist, lässt sich am besten durch Feedback von Arbeitskollegen erfassen. Da bereits Projekte mit dem Skeleton in Umsetzung sind und mehrere Arbeitskollegen des Prüflings mit dem Code arbeiten, wurde viel Feedback generiert. Dieses viel größtenteils positiv aus.

\subsection{Pflichtenheft}
\label{sec:Pflichtenheft}

Ein Pflichten- oder Lastenheft wurde für dieses Projekt nicht erstellt. 


% !TEX root = ../Projektdokumentation.tex
\section{Implementierungsphase} 
\label{sec:Implementierungsphase}

\subsection{Docker-Setup}
\label{sec:Docker-Setup}

Als Image für das aktuellste Typo3 Docker-Setup wurde das Image von Martin Helmich genutzt (https://github.com/martin-helmich/docker-typo3). Ein Docker-Image ist eine Datei, welche aus Anweisungen besteht, welche vollständige und ausführbare Version einer Anwendung erstellt. Diese Anweisungen werden in einer docker-compose.yml Daten gespeichert und dann mit dem Befehl docker-compose up ausgeführt. Nach dem Ausführen des Befehls, wird eine Typo3 Instanz erstellt und gestartet. Abhängig von dem angegebenen Port, kann dann das Typo3 über Localhost aufgerufen werden. Nach dem ausführen der Anweisungen, wird dort der Installations-Screen(siehe Abbildung) von Typo3 angezeigt. Nach dem Installieren, kann ausgewählt werden, ob das Typo3 mit einer leeren Startseite gestartet werden soll, oder weitere Konfigurationen installiert/ausgeführt werden sollen. Um das Skeleton möglichst minimal zu halten, wird das Typo3 mit einer leeren Startseite initialisiert.

\subsection{Ausspielen der Typo3 Daten}
\label{sec:Ausspielen der Typo3 Daten}

Damit die Typo3 Daten im JSON-Format ausgespielt werden, wird die Headless-Extension \linebreak (https://extensions.typo3.org/extension/headless) installiert. Danach muss die Extension zum Template hinzugefügt werden, damit Typo3 weiß, dass es das Template dementsprechend anpassen muss. Typo3 spielt nun die Daten im JSON-Format aus. Das JSON muss noch angepasst werden, damit eine Navigationsstruktur damit abgebildet werden kann. Dafür werden folgende Erweiterungen am Template vorgenommen.

\colorlet{punct}{red!60!black}
\definecolor{background}{HTML}{EEEEEE}
\definecolor{delim}{RGB}{20,105,176}
\colorlet{numb}{magenta!60!black}

\lstdefinelanguage{json}{
    basicstyle=\normalfont\ttfamily,
    numbers=left,
    numberstyle=\scriptsize,
    stepnumber=1,
    numbersep=8pt,
    showstringspaces=false,
    breaklines=true,
    frame=lines,
    backgroundcolor=\color{background},
    literate=
     *{0}{{{\color{numb}0}}}{1}
      {1}{{{\color{numb}1}}}{1}
      {2}{{{\color{numb}2}}}{1}
      {3}{{{\color{numb}3}}}{1}
      {4}{{{\color{numb}4}}}{1}
      {5}{{{\color{numb}5}}}{1}
      {6}{{{\color{numb}6}}}{1}
      {7}{{{\color{numb}7}}}{1}
      {8}{{{\color{numb}8}}}{1}
      {9}{{{\color{numb}9}}}{1}
      {:}{{{\color{punct}{:}}}}{1}
      {,}{{{\color{punct}{,}}}}{1}
      {\{}{{{\color{delim}{\{}}}}{1}
      {\}}{{{\color{delim}{\}}}}}{1}
      {[}{{{\color{delim}{[}}}}{1}
      {]}{{{\color{delim}{]}}}}{1},
}
\begin{lstlisting}[language=json,firstnumber=1]
lib.page = JSON
lib.page {
  fields {
    navigation {
      fields {
        main {
          dataProcessing {
            10 = FriendsOfTYPO3\Headless\DataProcessing\MenuProcessor
            10 {              
				       special = categories
				       special.value = 1
				       levels = 7
				       as = main
				       expandAll = 1
				       includeSpacer = 1
				       titleField = nav_title // title	
              	dataProcessing {
                	10 = FriendsOfTYPO3\Headless\DataProcessing\FilesProcessor
                	10 {
                  		references.fieldName = media
                  		as = media
                	}
              	}
                 overwriteMenuLevelConfig {
                  stdWrap.cObject {
                    100 = TEXT
                    100.field = uid
                    100.wrap = ,"uid":|
                  }
                }   
            }
          }
        }
      }      
    }
  }
}
\end{lstlisting}
Durch diese Anpassungen erhält das page Objekt die property navigation. Die Eigenschaft navigation hat wiederrum die Eigenschaft main. Mit der Hilfe des Menuprocessor von der Headless Extension, werden in die main Eigenschaft alle Seiten + Unterseiten geladen, welche der 1. Kategorie zugewiesen werden (special = categories \& special value = 1). Falls Bilder für die Seiten gepflegt sind, werden diese ebenfalls durch den FilesProcessor ausgespielt. Diese werde in dem Skeleton aber nicht verwendet. Dies dient nur zum Nutzen von potentiellen zukünftigen Projekten. Die Anpassungen können sowohl in der Skeleton-Extension stattfinden, als auch im Typo3 Backend.

\subsection{Skeleton-Extension / Layout}
\label{sec:Skeleton-Extension / Layout}
Typo3 ermöglicht es Entwicklern das Typo3 mit eigenen Extensions zu erweitern. Diese Extensions müssen einer bestimmten Struktur folgen, welche von Typo3 vorgegeben ist. Die genaue Form lässt sich in der Typo3 Dokumentation finden. \newline https://docs.typo3.org/m/typo3/reference-coreapi/main/en-us/ExtensionArchitecture/Index.html
\paragraph{Datenmodell:}

Alle Komponenten, die gepflegt werden können, sollen für Desktop, Tablet und Mobile verschiedene Breiten gepflegt bekommen können. Ein Text soll beispielsweise nur ein drittel der Bildschirmbreite auf Desktop haben, aber 50\% auf einem Tablet und die volle Breite auf einem Mobiltelefon. Dafür erhalten alle Elemente drei Felder in der Datenbank:
\begin{lstlisting}[language=json,firstnumber=1]
CREATE TABLE tt_content (
        tx_responsive_mobile int(11) DEFAULT '100' NOT NULL,
        tx_responsive_tablet int(11) DEFAULT '100' NOT NULL,
        tx_responsive_desktop int(11) DEFAULT '100' NOT NULL,
);
\end{lstlisting}
Anpassungen an bestehende Tabellen in der Datenbank werden in der ext\_tables.sql Datei vorgenommen. Diese Datei wird automatisch beim installieren von Extensions ausgelesen und die Anpassungen durchgeführt. Felder, welche als Input für Inhaltselemente dienen, werden in Typo3 in der tt\_content Tabelle gespeichert. Diese Felder müssen noch mit angemessenen Auswahlmöglichkeiten versehen werden und dann allen Elementen zugewiesen werden.
\paragraph{Backend:} Die Inputfelder werden als Selects dargestellt. Die Auswahlmöglichkeiten werden im Code vorgegeben.
\begin{lstlisting}[language=json,firstnumber=1]
"tx_responsive_mobile" => Array (
        "exclude" => 1,
        "label" => 'Mobile',
        "config" => Array (
            'type' => 'select',
            'renderType' => 'selectSingle',
            'items' => [
                ['25%','25'],
                ['33%','33'],
                ['50%','50'],
                ['66%','66'],
                ['75%','75'],
                ['100%','100']
            ],
            'default' => '100',
            'size' => 1,
            'maxitems' => 1,
        )
    ),
\end{lstlisting}
Analog zum oberen Code, werden für Tablet und Desktop ebenfalls Konfigurationen vorgenommen und im Array \$tempColumns gespeichert. Welche Eigenschaften Inputfelder brauchen und welche Optionen es bei den Inputfeldern gibt, lässt sich in der Typo3 Dokumentation(https://docs.typo3.org/m/typo3/\newline reference-tca/main/en-us/Columns/Index.html) nachlesen. Um die Inputfelder nun allen Elementen hinzuzufügen, wird die ExtensionManagementUtility Klasse aus dem Typo3 Core genutzt.
\begin{lstlisting}[language=json,firstnumber=1]
\TYPO3\CMS\Core\Utility\ExtensionManagementUtility::addTCAcolumns("tt_content",$tempColumns,1);
\TYPO3\CMS\Core\Utility\ExtensionManagementUtility::addToAllTCAtypes('tt_content','--div--;Responsive,tx_responsive_mobile,tx_responsive_tablet,tx_responsive_desktop','','after:addToAllTCAtypes');
\end{lstlisting}

Die erste Zeile Code speichert die Inputfelder ab, die zweite Zeile Code, fügt sie allen Inhaltselementen hinzu. --div--;Responsive sorgt dafür, dass diese in einem eigenen Tab mit dem Namen Responsive dargestellt werden. Danach wird aufgelistet werden Inputfelder hinzugefügt werden.

\paragraph{Ausspielen ans Frontend:}

Damit die gepflegten Daten auch ans Frontend ausgespielt werden, muss das Template von Typo3 um diese erweitert werden. Extensions haben dafür im Pfad \newline Configuration->TypoScript die Datei setup.typoscript. Hier können Anpassungen ans Template gepflegt werden, diese Anpassungen werden erst aktiv, wenn man die Extension dem Typo3 Template zuweist.

\begin{lstlisting}[language=json,firstnumber=1]
lib.appearance {
    fields {
        responsive_desktop = TEXT
        responsive_desktop {
            field = tx_responsive_desktop
        }
        responsive_tablet = TEXT ......
    }
}
\end{lstlisting}

Die gepflegten Daten werden nun in der Eigenschaft appearance ausgespielt und können von Nuxt3 ausgelesen werden.

\paragraph{Sections:}
Um die Elemente auf einer Seite in Bereiche auszuteilen, wurde ein komplett neues Inhaltselement erstellt. Dieses wurde als Section bezeichnet. Sections können Hintergrundfarben haben, wodurch Abschnitte im Frontend farblich voneinander getrennt dargestellt werden können. Dafür wurde ebenfalls in der tt\_content Tabelle ein Feld angelegt. In dem Fall mit dem Namen tx\_bal\_column\_color. Damit das neue Inhaltselement in Typo3 genutzt werden kann, muss die tt\_content.php Datei um folgenden Code erweitert werden. \Anhang{tt_content.php} \newline Die Datei lässt sich im Pfad Configuration->TCA->Overrides der Extension finden.

Damit das neue Element auch in der Auswahl von Elementen angezeigt wird, muss das Element noch in der ext\_localconf.php Datei im Root Verzeichnes der Extension hinzugefügt werden.\newline \Anhang{ext_localconf.php} \newline Übersetzungen für den Namen und die Beschreibung der Section wurden in einer .xlf Datei gespflegt.


\subsection{Implementierung der Benutzeroberfläche}
\label{sec:ImplementierungBenutzeroberflaeche}

\begin{itemize}
	\item Beschreibung der Implementierung der Benutzeroberfläche, falls dies separat zur Implementierung der Geschäftslogik erfolgt (\zB bei \acs{HTML}-Oberflächen und Stylesheets).
	\item \Ggfs Beschreibung des Corporate Designs und dessen Umsetzung in der Anwendung.
	\item Screenshots der Anwendung
\end{itemize}
\paragraph{Nuxt3 aufsetzen um DaisyUI installieren}
Um Nuxt3 nutzen zu können, muss eine aktuelle Version von Node.js auf der Betriebssystem installiert werden. Zum installieren von Nuxt3 und DaisyUI wurde der Node Package Manager genutzt.

\paragraph{Dynamische Seiten}
Alle Seiten der Nuxt3 Anwendungen befinden sich im pages Ordner. Da Nuxt3 nicht weiß, wie der Seitenbaum von Typo3 aussieht und dieser sich ständig ändert, können die einzelnen Seiten nicht per Hand eingepflegt werden. Damit Nuxt3 trotzdem weiß welcher Code bei den unterschiedlichen Seiten ausgeführt werden soll, gibt es eine Catch-all Route. Eine Datei/Seite die den Namen [...slug].vue erhält, fungiert als Catch-all Route. Dies bedeutet, dass alle Seiten die nicht gepflegt sind, durch diese Datei verarbeitet werden. Wenn aber beispielsweise eine Shop-Seite angelegt werde soll, die ihre Daten durch eine E-Commerce Platform bekommt. Dann kann im pages Ordner eine shop.vue Datei hinterlegt werden, welche dann für den Pfad beispielseite.de/shop die Daten darstellt. Analog kann dies für alle möglichen Seiten geschehen.

\paragraph{Verarbeitung der Typo3 Daten}

In der nuxt.config.ts Datei können generelle Konfigurationen am Nuxt3 vorgenommen werden. In ihr wurde die URL des Typo3 gespeichert. Dies hat den Vorteil, dass die URL in der Anwendung als Variable genutzt werden kann. Das heißt, dass bei späteren Entwicklungen lediglich die Variable einen anderen Wert bekommen muss und Nuxt automatisch alle Daten von einem anderen Typo3 bezieht.

\begin{lstlisting}[language=json,firstnumber=1]
export default defineNuxtConfig({
    runtimeConfig: {
        public: {
            typo3: 'http://localhost'
        }
    },
})
\end{lstlisting}

Die Variable wird darauf hin in der [...slug].vue Datei genutzt um die Daten vom Typo3 zu laden. Dies geschieht im <script setup> Tag. Nuxt3 Dateien können sowohl einen Script Tag haben, welcher Clientseitig ausgeführt wird, als auch ein Script Tag, welcher Serverseitig ausgeführt wird(hat den Zusatz setup). Da die ladezeit minimal gehalten werden soll, werden die Typo3 Daten hier Serverseitig geladen. Für \newline Entwicklungszwecke wird der Cache hier noch ausgeschaltet. Damit Änderungen im Typo3 direkt sichtbar werden. Wenn das Projekt in Production geht, sollte die Zeile {initialCache: false} entfernt werden. Dadurch Speichert Nuxt3 die Daten im Cache und die Seite wird schnell geladen. Je nach Website kann programmiert werden, dass sich der Cache häufiger oder weniger häufiger leert.

\begin{lstlisting}[language=json,firstnumber=1]
  const runtimeConfig = useRuntimeConfig()
  const route = useRoute();
  const { data } = await useAsyncData(
      "pageData",
    () => $fetch(runtimeConfig.typo3 + route.fullPath),
    {initialCache: false}
  )
\end{lstlisting}

Im route Objekt ist der Pfad der aktuellen Website gespeichert. route.fullPath würde bei der Seite beispielseite.de/beispiel1 dem Wert beispiel1 entsprechen.
Die Daten aus dem Typo3(hier das data Objekt) wurden danach weiter verarbeitet. Der erste Schritt der Datenverarbeitung war es die Daten in Kategorien zu zerlegen. Dafür wurden vier Kategorien/Variablen angelegt. Die Variable breadcrumbs bekam die Daten bezüglich der Breadcrumbs der Seite. Die Variable Content bekam alle Daten bezüglich der Inhaltselemente auf der bestimmten Seite. Die Variable mainNavigation bekam alle Daten der Hauptnavigation. Als letztes bekam die Variable metaData alle Metadaten der Website. Die Breadcrumbs wurden bereits automatisch durch die Headless Extension generiert und sind im Datenobjekt vorhanden. Das gleiche gilt für die Metadaten und dem Großteil der Inhaltselemente. Die Navigationsdaten finden sich in der Eigenschaft main, welche zum Objekt navigation gehört, welches wiederum zum Objekt page gehört. Diese Logik wurde in \ref{sec:Ausspielen der Typo3 Daten} angelegt.

\begin{lstlisting}[language=json,firstnumber=1]
  let breadcrumbs = data.value.breadcrumbs;
  let content = data.value.content.colPos0;
  let mainNavigation = data.value.page.navigation.main;
  let metaData = data.value.meta;
\end{lstlisting}

Mit der Nuxt3 Funktion useHead() können die Metadaten an die Website weiter gegeben werden. Für die Weiterverarbeitung der Breadcrumbs, die Navigation und den Content wurden eigene Komponenten erstellt. Diese erhalten die Variablen als Eigenschaften/Properties/Props zugewiesen. Da Nuxt3 einen automatischen Import von Komponenten hat, müssen diese lediglich im Template gepflegt werden und nicht zusätzlich importiert werden. Dafür wurden die Komponenten im components Ordner gepflegt, damit Nuxt3 diese automatisch importieren kann. Den vollständigen Code der [...slug].vue finden Sie unter \Anhang{[...slug].vue}.

\paragraph{Finale Anwendung}
Screenshots der finalen Anwendung inklusive der gepflegten Dummy-Inhalte der Typo3 Seite befinden sich im \Anhang{Screenshots}.


\subsection{Implementierung der Geschäftslogik}
\label{sec:ImplementierungGeschaeftslogik}

\begin{itemize}
	\item Beschreibung des Vorgehens bei der Umsetzung/Programmierung der entworfenen Anwendung.
	\item \Ggfs interessante Funktionen/Algorithmen im Detail vorstellen, verwendete Entwurfsmuster zeigen.
	\item Quelltextbeispiele zeigen.
	\item Hinweis: Wie in Kapitel~\ref{sec:Einleitung}: \nameref{sec:Einleitung} zitiert, wird nicht ein lauffähiges Programm bewertet, sondern die Projektdurchführung. Dennoch würde ich immer Quelltextausschnitte zeigen, da sonst Zweifel an der tatsächlichen Leistung des Prüflings aufkommen können.
\end{itemize}

\paragraph{Beispiel}
Die Klasse \texttt{Com\-par\-ed\-Na\-tu\-ral\-Mo\-dule\-In\-for\-ma\-tion} findet sich im \Anhang{app:CNMI}.  

% !TEX root = ../Projektdokumentation.tex
\section{Abnahmephase} 
\label{sec:Abnahmephase}
Getestet wurde die Anwendung, indem sie Entwicklern gegeben wurde und diese damit probiert haben, Seiten zu Entwickeln. Diese 'Tests' verliefen sehr gut. Klassische Tests(Unit, Integration, etc.) werden je nach potentiell damit zu entwickelnden Projekten nachträglich erst eingefügt. Eine offizielle Abnahme im klassischen Sinne gab es in diesem Projekt nicht. Das lag nicht daran, dass das Projekt schlecht war, sondern an der Natur des Projekts und der Technologien. Das Projekt sollte dazu genutzt werden, neue Projekte zu entwickeln, welche dann von Kunden abgenommen werden. Das Projekt wurde in dem Sinne abgenommen, dass die Arbeitskollegen des Prüfungsbewerbers es zum aktuellen Zeitpunkt nutzen und damit neue Projekte umsetzen. Eine richtige Abnahme / Übergabe an einem Kunden gab es aber nicht. Die Anwendung wird wahrscheinlich abhängig vom Feedback der Arbeitskollegen, auch in der Zukunft weitere Iterationen durchleben.

% !TEX root = ../Projektdokumentation.tex
\section{Einführungsphase}
\label{sec:Einfuehrungsphase}

Für das Nutzen des Projekts benötigen die Arbeitskollegen vom Prüfungsbewerber zwei Sachen. Als erstes brauchen sie alle Voraussetzungen um Nuxt 3 auf ihrem System laufen zu lassen. Diese finden sich in der Nuxt 3 Dokumentation. Danach kann das repository von Github geklont werden und das Skeleton Clientseitig genutzt werden.

Als zweites wird ein Typo3 gebraucht. Dies kann ein bestehendes Typo3 auf einem Server sein, welches auf headless eingestellt wurde. Dies kann aber auch ein Typo3 sein, welches lokal in einem Docker-Container läuft.

Um das finale Projekt zu deployen, welches aus dem Skeleton entwickelt wurde, bedarf es nur zwei Dinge. Der Nuxt 3 Code muss den Build-Prozess durchlaufen. Dafür wird im entsprechenden Verzeichnis einfach der Befehl nuxt build im Terminal ausgeführt. Dadurch entsteht ein output Verzeichnis. Im output Verzeichnis wird durch den Befehl node .output/server/index.mjs der Nuxt 3 Server gestartet. 

Typo3 kann auf einem Server installiert werden, welcher PHP nutzen kann. Genauere Anforderungen finden sich in der Typo3 Dokumentation. Oder es kann ebenfalls in einem Docker-Container laufen. Entsprechende Images finden sich im Internet. Wenn Änderungen am Typo3 werden im Normalfall durch das installieren von Extensions durchgeführt, was im Backend vom Typo3 geschieht. Dadurch muss ein einmalig laufendes Typo3 nicht komplett neu gestartet werde, wenn Anpassungen gemacht werden. Wenn Anpassungen am Nuxt 3 erfolgen, muss dieses neu gebaut werden(Build-Prozess) und der neu gebaute Code
% !TEX root = ../Projektdokumentation.tex
\section{Dokumentation}
\label{sec:Dokumentation}

Wie in den Unterpunkten vom Anhang \ref{app:Doc} Entwicklerdokumentation zu finden ist, wurde für dieses Projekt eine Entwicklerdokumentation angelegt. Es wurde keine Kundendokumentation angelegt. Dies hat den Grund, dass die Abnehmer des Projekts Entwickler der bits \& likes GmbH sind. Aus den in der Zukunft erstellten Projekten sollten aber Kundendokumentationen erstellt werden. Da diese aber komplett unterschiedliche Designs und neue Komponenten haben werden, ergibt es kein Sinn in dem Kontext des Projektes eine zu erstellen. Der Großteil müsste umgeschrieben werden, wenn Projekte mit dem Skeleton umgesetzt werden. Für die Entwickler wurde dokumentiert, wie ein Typo3 in Docker gestartet werden kann und wie das Typo3 mit Nuxt 3 interagiert. Die Entwickler sind damit dann in der Lage das Projekt zu erweitern und komplexere Designs umzusetzen.
% !TEX root = ../Projektdokumentation.tex
\section{Fazit} 
\label{sec:Fazit}

\subsection{Soll-/Ist-Vergleich}
\label{sec:SollIstVergleich}

\begin{itemize}
	\item Wurde das Projektziel erreicht und wenn nein, warum nicht?
	\item Ist der Auftraggeber mit dem Projektergebnis zufrieden und wenn nein, warum nicht?
	\item Wurde die Projektplanung (Zeit, Kosten, Personal, Sachmittel) eingehalten oder haben sich Abweichungen ergeben und wenn ja, warum?
	\item Hinweis: Die Projektplanung muss nicht strikt eingehalten werden. Vielmehr sind Abweichungen sogar als normal anzusehen. Sie müssen nur vernünftig begründet werden (\zB durch Änderungen an den Anforderungen, unter-/überschätzter Aufwand).
\end{itemize}

Im Groben wurden alle Projektziele erreicht. Es existiert jetzt eine Nuxt 3 Anwendung, welche sich mit einem Headless-Typo3 verbindet und dessen content ausspielt. Das einzige, was leider nicht mehr in der Projektzeit umzusetzen war, ist die Dokumentation für Entwickler. Als temporärer Ersatz kann diese Dokumentation dienen. Einige Projektphasen haben etwas mehr oder weniger Zeit gebraucht, es ist aber alles in einem zu erwarteten Rahmen geblieben(< +-2 Stunden).

\tabelle{Soll-/Ist-Vergleich}{tab:Vergleich}{Zeitnachher.tex}


\subsection{Lessons Learned}
\label{sec:LessonsLearned}

\begin{itemize}
	\item Was hat der Prüfling bei der Durchführung des Projekts gelernt (\zB Zeitplanung, Vorteile der eingesetzten Frameworks, Änderungen der Anforderungen)?
\end{itemize}

\subsection{Ausblick}
\label{sec:Ausblick}

\begin{itemize}
	\item Wie wird sich das Projekt in Zukunft weiterentwickeln (\zB geplante Erweiterungen)?
\end{itemize}


% Literatur ------------------------------------------------------------------
\clearpage
\renewcommand{\refname}{Literaturverzeichnis}
\bibliography{Bibliographie}
\bibliographystyle{Allgemein/natdin} % DIN-Stil des Literaturverzeichnisses
\cleardoublepage
% !TEX root = Projektdokumentation.tex
\clearpage
\addsec{Eidesstattliche Erklärung}

% Hinweis: die eidesstattliche Erklärung ist ggfs. an die Richtlinie der IHK anzupassen

Ich, \autorName, versichere hiermit, dass ich meine \textbf{\betreff} mit dem
Thema
\begin{quote}
\textit{\kompletterTitel}
\end{quote}
selbständig verfasst und keine anderen als die angegebenen Quellen und Hilfsmittel benutzt habe,
wobei ich alle wörtlichen und sinngemäßen Zitate als solche gekennzeichnet habe. Die Arbeit
wurde bisher keiner anderen Prüfungsbehörde vorgelegt und auch nicht veröffentlicht.\\[6ex]

\abgabeOrt, den \abgabeTermin


\rule[-0.2cm]{5.5cm}{0.5pt}

\textsc{\autorName}

\cleardoublepage
\phantomsection
\thispagestyle{empty}
\pdfbookmark[1]{Eidesstattliche Erklärung}{ihkdeckblatt}
\includegraphicsKeepAspectRatio{DeckblattIHK}{1}
\cleardoublepage

% Anhang ---------------------------------------------------------------------
\clearpage
\appendix
\pagenumbering{roman}
% !TEX root = Projektdokumentation.tex
\section{Anhang}

\begin{figure}[htb]
\centering
\includegraphicsKeepAspectRatio{FlussDiagramm}{1}
\caption{FlussDiagramm Nuxt 3 Anwendung}
\label{fig:FlussDiagramm}
\end{figure}

\pagebreak

\subsection{tt\_content.php}
\label{tt_content.php}

\begin{lstlisting}[language=json,firstnumber=1]
<?php declare(strict_types=1);

defined('TYPO3_MODE') || die();

(static function () {
    \TYPO3\CMS\Core\Utility\ExtensionManagementUtility::addPlugin(
        array(
            'LLL:EXT:bal_skeleton/Resources/Private/Language/Tca.xlf:bal_column.wizard.title',
            'tx_bal_column',
            'EXT:bal_skeleton/Resources/Public/Icons/ContentElements/stage.png'
        ),
        'CType',
        'bal_skeleton'
    );
    $temporaryColumn = array(
        'tx_bal_column_color' => array (
            'exclude' => 1,
            'label' => 'LLL:EXT:bal_skeleton/Resources/Private/Language/Tca.xlf:bal_column.color.title',
            'config' => array (
                'type' => 'input',
                'renderType' => 'colorpicker',
                'size' => 10,
            )
        ),
    );
    \TYPO3\CMS\Core\Utility\ExtensionManagementUtility::addTCAcolumns(
        'tt_content',
        $temporaryColumn
    );
    $GLOBALS['TCA']['tt_content']['types']['tx_bal_column'] = array(
        'showitem' => '
            --palette--;LLL:EXT:frontend/Resources/Private/Language/locallang_ttc.xml:palette.general;general,
            tx_bal_column_color,
    ');
})();
\end{lstlisting}

\subsection{ext\_localconf.php}
\label{ext_localconf.php}

\begin{lstlisting}[language=json,firstnumber=1]
<?php
defined('TYPO3_MODE') || die('Access denied.');

use TYPO3\CMS\Extbase\Utility\ExtensionUtility;

call_user_func(
    function () {
        // wizards
        \TYPO3\CMS\Core\Utility\ExtensionManagementUtility::addPageTSConfig(
            'mod.wizards.newContentElement.wizardItems.common {
           elements {
              tx_bal_column {
                 iconIdentifier = bal_column
                 title = LLL:EXT:bal_skeleton/Resources/Private/Language/Tca.xlf:bal_column.wizard.title
                 description = LLL:EXT:bal_skeleton/Resources/Private/Language/Tca.xlf:bal_column.wizard.description
                 tt_content_defValues {
                    CType = tx_bal_column
                 }
              }
           }
           show := addToList(tx_bal_column)
        }'
        );

        $iconRegistry = \TYPO3\CMS\Core\Utility\GeneralUtility::makeInstance(\TYPO3\CMS\Core\Imaging\IconRegistry::class);

        $iconRegistry->registerIcon(
            'bal_column',
            \TYPO3\CMS\Core\Imaging\IconProvider\SvgIconProvider::class,
            ['source' => 'EXT:bal_skeleton/Resources/Public/Icons/ContentElements/stage.svg']
        );
   }
);
\end{lstlisting}

\subsection{[...slug].vue}
\label{[...slug].vue}

\begin{lstlisting}[language=json,firstnumber=1]
<script setup lang="ts">
  const runtimeConfig = useRuntimeConfig()
  const route = useRoute();
  const { data } = await useAsyncData(
      "pageData",
    () => $fetch(runtimeConfig.typo3 + route.fullPath),
    {initialCache: false}
  )
  let breadcrumbs = data.value.breadcrumbs;
  let content = data.value.content.colPos0;
  let mainNavigation = data.value.page.navigation.main;
  let metaData = data.value.meta;

  useHead({
    title: metaData.title,
    meta: [
      { name: 'description', content: metaData.description ? metaData.description : '' }
    ]
  })

</script>
<template>
     <div>
          <template v-if="mainNavigation"><Navigation :typo3-navigation="mainNavigation"></Navigation></template>
          <div class="container mx-auto xl">
            <template v-if="breadcrumbs"><Breadcrumbs :typo3breadcrumbs="breadcrumbs"></Breadcrumbs></template>
            <template v-if="content"><TypoView :typo3-content="content"></TypoView></template>
            <h3 v-else>Noch kein Content gepflegt!</h3>
          </div>
     </div>
</template>
\end{lstlisting}

\pagebreak

\subsection{navigation.ts}
\label{navigation.ts}

\begin{lstlisting}[language=json,firstnumber=1]
export const useNavigation = () => {
    const fullNavList = useState('fullNavList', () => []);
    const navigationOpened = useState('navigationOpened', () => false);
    const currentNavList = useState('currentNavList', () => []);
    const previousNavItem:{value: navigationItem | Record<string, never>} = useState('previousNavItem');
    const router = useRouter();
    
    const open = (navItem: navigationItem) => {
        if (!navItem.hasSubpages) {
            router.push(navItem.link)
        } else {
            previousNavItem.value = navItem;
            currentNavList.value = navItem.children;
            navigationOpened.value = true;  
        } 
    }
    const close = () => {
        navigationOpened.value = false;
    }
    const navigate = (navItem: navigationItem) => {
        if (!navItem.hasSubpages) {
            close();
            router.push(navItem.link)
        }
        else {
            previousNavItem.value = navItem;
            currentNavList.value = navItem.children;
        }
    }
    const back = () => {
        const parent = findParent(fullNavList.value, previousNavItem.value.uid);
        if (parent && parent.children) {
            previousNavItem.value = parent;
            currentNavList.value = parent.children;
        } else {
            previousNavItem.value = {};
            currentNavList.value = fullNavList.value;
        }
    }
    const findParent = (nav: navigationItem[], id: number, potentialParentNav?: navigationItem) => {
        let parent: navigationItem;
        for (let index = 0; index < nav.length; index++) {
            const item = nav[index];
            if (item.uid === id) {
                parent = potentialParentNav;
                break;
            } 
            else {
                if (item.children && item.children.length > 0) {
                        parent = findParent(item.children, id, item);
                        if (parent) {
                            break;
                        }
                    }
                }
        }
        return parent;
    }
    return {fullNavList, navigationOpened, open, close, currentNavList, previousNavItem, navigate, back};
}
\end{lstlisting}

\subsection{SlideMenu.vue}
\label{SlideMenu.vue}

\begin{lstlisting}[language=json,firstnumber=1]
<script setup>
    const {fullNavList, navigationOpened, close, currentNavList, previousNavItem, navigate, back} = useNavigation();
</script>
<template>
    <div class="slidemenu-container" :class="{'active' : navigationOpened}">
        <div class="flex">
            <div class="text-xl font-bold cursor-pointer" v-if="fullNavList != currentNavList" @click="back()">back</div>
            <div class="text-xl font-bold pl-6 grow" v-if="previousNavItem">
                <NuxtLink :to="previousNavItem.link" @click="close()">{{previousNavItem.title}}</NuxtLink>
            </div>
            <div class="text-xl font-bold cursor-pointer" @click="close()">X</div>
        </div>
        <hr>
        <div class="navbar">
            <div class="btn btn-ghost normal-case text-xl" v-for="navItem in currentNavList" :key="navItem.uid" @click="navigate(navItem)">
                {{navItem.title}}
            </div>
        </div>
    </div>
</template>
<style>
    .slidemenu-container {
        position: absolute;
        height: 100vh;
        width: 0px;
        transition: width 0.5s;
        background: white;
        overflow-x: hidden;
        z-index: 100;
    }
    .slidemenu-container.active {
        width: 400px;
    }
</style>
\end{lstlisting}
\clearpage
\input{Anhang/AnhangMenü.tex}
\subsection{Screenshots der Anwendung}
\label{Screenshots}
\begin{figure}[htb]
\centering
\includegraphicsKeepAspectRatio{tagliste.pdf}{1}
\caption{Anzeige und Filterung der Module nach Tags}
\end{figure}
\clearpage
\begin{figure}[htb]
\centering
\includegraphicsKeepAspectRatio{modulliste.pdf}{1}
\caption{Liste der Module mit Filtermöglichkeiten}
\end{figure}
\clearpage

\subsection{Entwicklerdokumentation}
\label{app:Doc}

\paragraph{lokales Setup}

Um das Skeleton lokal zu starten, müssen zwei Dinge geschehen. Es muss das Typo3 lokal gestartet werden, alternativ kann auch ein vorhandenes Typo3 auf einem Server angesteuert werden. Als zweites muss Nuxt 3 gestartet werden und mit dem Typo3 Verbunden werden. 

Um Typo3 lokal zu starten, wurde ein Docker-Image erstellt, welches sich auf Github runterladen lässt. Mit dem Befehl docker-compose up lässt sich dieses starten. Je nachdem welcher Port im Image angegeben wurde kann das Typo3 jetzt über localhost angesteuert werden.

Um Nuxt 3 lokal zu starten, bedarf es kein Docker. Docker kann aber auch genutzt werden, wenn dies erwünscht ist. Dafür muss sich das aktuelle Github Repository heruntergeladen werden. Mit dem Befehl npm i(angenommen node.js/npm ist installiert), werden die benötigten Pakete heruntergeladen um Nuxt 3 zu starten. Mit dem Befehl npm run dev kann nun Nuxt 3 gestartet werden.

\paragraph{Typo3 bereits machen}

Um die Ausgabe von Typo3 in ein JSON-Format umzuwandeln, muss noch die Headless-Extension installiert werden. Um die ausgespielten Daten in die Struktur zu bringen, die der Nuxt 3 Client erwartet, muss ebenfalls die Skeleton-Extension installiert werden.

\paragraph{Verbindung zwischen Typo3 und Nuxt 3}

Um den Nuxt 3 Clienten mit dem lokalen Typo3 zu verbinden, muss in der nuxt.config.ts noch die entsprechende URL gepflegt werden. Dies geschieht in dem runtimeConfig.public Objekt. Dort ist eine Eigenschaft mit dem Namen typo3 gepflegt. In ihr kann die Typo3 URL ausgetauscht werden. Die Daten des Typo3 werden in der [...slug].vue Datei im pages Verzeichnis abgerufen.

\paragraph{Verarbeitung der Typo3 Daten}

Wie die Daten von dem Nuxt 3 verarbeitet werden und die entsprechenden Elemente / Navigation dargestellt wird, lässt sich an dem Flussdiagramm erkennen.




\end{document}
